% !TEX program = lualatex
\documentclass[aspectratio=43,professionalfonts,handout]{beamer}
\usepackage{beamer_template}

\newcommand{\LectureNum}{15}
\newcommand{\LectureTitle}{一般角と三角関数}
\newcommand{\TextPages}{pp.\,140-142}
\newcommand{\NextTopic}{弧度法}
\newcommand{\NextPages}{pp.\,143-145}
\newcommand{\CourseName}{基礎数学B}
\renewcommand{\TermName}{後期}

\begin{document}

% -----------------------------------------------
% スライド1: 表紙＋目標＋用語・定義
% -----------------------------------------------
\begin{frame}{\CourseName\quad 第\LectureNum 回「\LectureTitle 」}
  {\small \TermName\quad \TeacherName\quad ／\quad 教科書 \TextPages}

  \begin{block}{用語・定義}
  \begin{description}[labelwidth=5.5em]
    \item[\kw{一般角：$360° \times n + \theta$ で表す}] $360° \times n + \theta$（$n$ は整数）で表す角
    \item[\kw{動径と象限}] 角 $\theta$ の動径がある象限で三角関数の符号が決まる
    \item[\kw{単位円上の定義：$\sin\theta = y$、$\cos\theta = x$、$\tan\theta = y/x$}] 単位円上の点 $(\cos\theta, \sin\theta)$；$\tan\theta = \sin\theta / \cos\theta$
  \end{description}
  \end{block}
\end{frame}

% -----------------------------------------------
% スライド2: 公式・性質
% -----------------------------------------------
\begin{frame}{公式・性質}
  \begin{block}{}
  \begin{itemize}
    \item $\sin$・$\cos$ の周期：$2\pi$（$360°$）
    \item $\tan$ の周期：$\pi$（$180°$）
    \item 各象限での符号
  \end{itemize}
  \end{block}
\end{frame}

% -----------------------------------------------
% スライド3: 例1→問1
% -----------------------------------------------
\begin{frame}{例1→問1}
  \begin{exampleblock}{例1) p.148（動径の位置・象限）}
    （板書で解説）
  \end{exampleblock}

  \vfill

  \begin{block}{問1) p.148\quad 自分で解いてみよう}
    （ノートに解くこと）
  \end{block}
\end{frame}

% -----------------------------------------------
% スライド4: 例2→問2
% -----------------------------------------------
\begin{frame}{例2→問2}
  \begin{exampleblock}{例2) p.149（一般角の三角関数の値）}
    （板書で解説）
  \end{exampleblock}

  \vfill

  \begin{block}{問2) p.149\quad 自分で解いてみよう}
    （ノートに解くこと）
  \end{block}
\end{frame}

% -----------------------------------------------
% まとめ
% -----------------------------------------------
\begin{frame}{まとめ}
  \begin{itemize}
    \item 一般角 $\theta + 360° n$ で1回転以上や負の角を統一的に扱う
    \item 単位円：$\cos\theta = x$，$\sin\theta = y$，$\tan\theta = y/x$
    \item 象限によるサインの符号（All/Sin/Tan/Cos の順）を確認する
  \end{itemize}

  \vfill
  {\small 基本問題：285, 286, 287}
\end{frame}

\end{document}
